\documentclass[aps,prd,twocolumn,10pt,nofootinbib,superscriptaddress]{revtex4-2}
\usepackage{amsmath,amssymb,graphicx}
\usepackage{booktabs}
\usepackage[colorlinks=true,linkcolor=blue,citecolor=blue,urlcolor=blue]{hyperref}

\begin{document}

\title{Gravitational lensing is not independent evidence for dark matter:\\
the no-slip identity in a cosmological-constant acceleration-scale framework}
\author{C.\ Zimmerman}
\affiliation{Independent}
\date{June 2026}

\begin{abstract}
Gravitational lensing is routinely presented as direct, model-independent evidence for particle dark matter---the
Bullet cluster's mass--gas offset, cluster strong$+$weak-lensing masses, the galaxy--galaxy-lensing radial
acceleration relation, and strong-lens Einstein radii. We show that within a framework in which the modified-dynamics
acceleration scale is set by the cosmological constant, $a_0=c^2\sqrt{\Lambda/32\pi}=9.36\times10^{-11}\,\mathrm{m\,
s^{-2}}$, realised covariantly by the Aether-Scalar-Tensor (AeST) theory, this evidentiary status collapses. AeST has
no gravitational slip: the two weak-field potentials are equal, $\Phi=\Psi$, so the lensing mass equals the dynamical
mass for any source. Four of the standard lensing-for-dark-matter arguments therefore reduce to a single statement---
that clusters carry a residual between dynamical and baryonic mass---rather than constituting four independent proofs.
We verify the identity from the AeST field equations and the primary literature, audit the acceleration-scale footing
of the measurements (cluster lensing masses are pure general-relativity$+$dark-matter deductions with no $a_0$; the
framework uses its own least-favourable $a_0$), and report the standing both ways: the galaxy--galaxy-lensing relation
passes; the cluster residual ($\eta\!\simeq\!2.3$) and the cosmic-microwave-background third peak remain genuinely
independent, conceded losses; the once-headline $8.6$--$9.2\sigma$ weak-lensing morphology split is now contested by
a re-analysis of the same data. Lensing, in this framework, is not a dark-matter detector but a second view of the
one shared cluster residual.
\end{abstract}

\maketitle

\section{The claim under examination}
Among the empirical pillars of cold dark matter, gravitational lensing occupies a special rhetorical place: it is
said to weigh mass directly, through the deflection of light, independently of any assumption about dynamical
equilibrium. The canonical statements are (a) the Bullet cluster 1E\,0657$-$558, whose weak- and strong-lensing
convergence peaks are spatially offset from the dominant baryonic component (the X-ray gas) and track the
collisionless galaxies~\cite{Clowe2006}; (b) cluster strong$+$weak-lensing masses that exceed the baryonic mass by a
factor $\sim2$; (c) the galaxy--galaxy-lensing radial acceleration relation (RAR), which extends the rotation-curve
RAR to accelerations below those reached kinematically~\cite{Brouwer2021}; and (d) strong-lens Einstein radii and
time delays. Each is presented as an \emph{independent} confirmation that gravitating mass exceeds visible mass.

The premise behind ``independent'' is that lensing and dynamics probe different things. In general relativity with a
dark halo they do not: both follow the same total mass. But for \emph{modified gravity} the premise is non-trivial,
because the original non-relativistic Milgromian dynamics (MOND~\cite{Milgrom1983}) famously \emph{under}-predicted
lensing---a real historical failure. The question this paper settles, for the present framework, is whether its
\emph{relativistic} completion repairs that failure in a way that also dissolves the ``independent evidence'' claim.

\section{The framework and its relativistic lensing}
The framework posits that the MOND acceleration scale is not a free constant but is fixed by the cosmological
constant,
\begin{equation}
a_0=c^2\sqrt{\Lambda/32\pi}=\tfrac{c}{2}\sqrt{G\rho_{\rm DE}}=cH_\Lambda/Z,\quad Z=\sqrt{\tfrac{32\pi}{3}},
\label{eq:a0}
\end{equation}
giving $a_0=9.36\times10^{-11}\,\mathrm{m\,s^{-2}}$ for the measured $\Lambda$; the coefficient $Z=5.79$ is an
unforced posit, flagged here and not asserted as derived. The non-relativistic dynamics follow the de Sitter--Unruh
interpolation $g_{\rm obs}=\sqrt{g_{\rm bar}^2+g_{\rm bar}a_0}$. The covariant realisation is the Aether-Scalar-Tensor
theory (AeST)~\cite{Skordis2021}, a Lorentz-violating scalar--vector--tensor theory whose unit-timelike vector field
is the device that fixes the lensing.

\paragraph{No gravitational slip.} Write the perturbed Friedmann--Robertson--Walker metric in conformal Newtonian
gauge as $ds^2=-(1+2\Phi)dt^2+(1-2\Psi)d\mathbf{x}^2$. Dynamics (non-relativistic test bodies) respond to $\Phi$;
light deflection responds to $\Phi+\Psi$. In general MOND-like theories these can differ (a ``slip'' $\Phi\neq\Psi$),
which is exactly how the old non-relativistic MOND under-lensed. AeST is constructed so that they do not. The
quasistatic solutions of AeST~\cite{Verwayen2024} state it directly: ``the two standard weak-field metric potentials
$\Phi$ and $\Psi$ are equal $\ldots$ this ensures that the lensing mass is equal to the dynamical mass $\ldots$ for
any source configuration where the MOND limit of AeST successfully replaces dark matter, the gravitational lensing
signal will also be as if dark matter is present.'' Equivalently, a single effective potential enters both metric
components~\cite{Mistele2023}; the unit-timelike vector sources $\Psi$ alongside $\Phi$, the Sanders mechanism. Thus
\begin{equation}
\Phi=\Psi\ \Longrightarrow\ M_{\rm lens}=M_{\rm dyn}\quad\text{(for any source).}
\label{eq:noslip}
\end{equation}
This is a field-equation identity, not a fit. Its only caveat is in the far field: beyond a transition radius the AeST
potentials enter an oscillatory regime~\cite{Verwayen2024}, and weak-lensing tests at the largest radii constrain the
amplitude~\cite{Mistele2023}; neither introduces a slip at the radii relevant below.

\section{The reframe}
Equation~\eqref{eq:noslip} removes the logical independence of the lensing arguments. If lensing mass equals dynamical
mass identically, then a lensing measurement that finds ``more mass than baryons'' finds \emph{the same} excess the
dynamics already show---not a second, corroborating excess. Table~\ref{tab:reframe} sorts the standard arguments
accordingly.

\begin{table}[t]
\caption{The lensing-for-dark-matter arguments under the no-slip identity. ``Reduces'' $=$ not independent of the
cluster dynamical residual; ``independent'' $=$ untouched by Eq.~\eqref{eq:noslip}.}
\label{tab:reframe}
\begin{ruledtabular}
\begin{tabular}{lll}
Argument & Status & Note \\
\colrule
Bullet mass--gas offset & reduces & collisionless residual \\
Cluster SL$+$WL $\sim2\times$ & reduces & $=$ dynamical $\eta$ \\
Galaxy--galaxy lensing RAR & reduces & passes (below) \\
Strong-lens Einstein radii & reduces & same potential \\
CMB third peak / $P(k)$ & independent & loss, $a_0$-free \\
Lensed-quasar substructure & independent & open \\
\end{tabular}
\end{ruledtabular}
\end{table}

Four arguments collapse into one: \emph{do clusters carry a dynamical residual?} They do---$\eta(R_{500})\!\simeq\!2.3$
on the eROSITA eRASS1 sample of 9830 clusters~\cite{Bulbul2024}---and that residual is the well-known, shared
liability of all modified-dynamics theories, not a new lensing discovery. Two arguments survive as genuinely
independent: the matter power spectrum encoded in the cosmic-microwave-background acoustic peaks (where $a_0$ is
absent from the linear physics, so the third-peak height is a real, conceded loss); and the abundance of
$10^7$--$10^9\,M_\odot$ substructure inferred from lensed-quasar flux-ratio anomalies, which probes clumpiness rather
than the smooth potential and is open in both directions.

\section{The acceleration-scale footing}
A recurring concern is whether the comparison silently uses a non-framework $a_0$. We audited it. Cluster
weak-lensing masses (eRASS1, DES-Y3, HSC) are computed in pure general relativity with a dark halo and contain
\emph{no} $a_0$: the inferred ``mass needed'' is simply the lensing-deduced enclosed mass exceeding the baryons. The
framework comparison then uses Eq.~\eqref{eq:a0}, $a_0=9.36\times10^{-11}$, which is the \emph{least}-favourable
choice, because in the deep regime $\eta\!\propto\!a_0^{-1/2}$: the canonical MOND value $1.2\times10^{-10}$ would make
the residual appear $\sim13\%$ \emph{smaller}. On the eRASS1 data we find $\eta(R_{500})=2.33$ on the framework's own
interpolation, versus $1.92$ at the canonical value---so the literature, which uses the canonical value, if anything
\emph{understates} the framework's residual. The galaxy--galaxy-lensing estimator is itself $a_0$-free ($g_{\rm
obs}\propto\Delta\Sigma$), and the morphology split below is $a_0$- and interpolation-independent by construction.
There is no footing error inflating any lensing loss; the only mis-footings we found ran the other way (a
galaxy-lensing amplitude and a cluster residual once displayed at the canonical $a_0$, flattering the framework), and
correcting them makes its numbers slightly worse---which we adopt.

\section{The Bullet cluster}
The Bullet is the iconic case. Under Eq.~\eqref{eq:noslip} its lensing peaks lie where the total gravitating mass is,
which in modified gravity is sourced by all matter; the peaks sitting on the collisionless galaxies rather than the
shocked gas requires that the cluster residual be \emph{collisionless} and galaxy-tracking---cold gas, stellar
remnants, or relic neutrinos---which is the standard residual, offset from the gas because the gas alone is
collisional. Recent consistent-QUMOND lens modelling reproduces the convergence-on-galaxies morphology with a residual
of order the cluster mass~\cite{Cha2026}; we treat this as imported rather than derived. The Bullet is thus a mild,
shared liability, not a clean falsification of modified gravity---and its $\sim4000\,\mathrm{km\,s^{-1}}$ collision
velocity remains a long-standing tension for $\Lambda$CDM pairwise velocities that survives the 2025 reclassification
of the system as a minor merger~\cite{JWSTbullet}. The complementary external-field prediction for the
dark-matter-deficient galaxy NGC\,1052-DF2 ($\sigma\!\sim\!14\,\mathrm{km\,s^{-1}}$) is a clean modified-dynamics
success.

\section{The morphology split}
One lensing tension does \emph{not} reduce: the galaxy--galaxy-lensing RAR split by morphology. In the released KiDS
catalogue, early-type galaxies lie $+0.26$ dex above late types at fixed baryonic acceleration---an $8.6$--$9.2\sigma$
contrast~\cite{Brouwer2021}. It is $a_0$- and interpolation-independent: a property-blind force law predicts a single
relation, and we confirm the early-minus-late model difference cancels identically for every $a_0$ and interpolation.
An intrinsic split of this size would falsify the law. It is, however, now \emph{contested}: a re-analysis of the same
data by the RAR's own authors, applying a stricter isolation cut, finds the early/late deviation falling from
$\sim7.5\sigma$ to $\sim0$ between isolation radii of $2.7$ and $4.0$ Mpc (the original used $3$ Mpc), attributing the
residual to early-type two-halo and satellite contamination from morphology--density clustering plus a mild
stellar-initial-mass-function offset~\cite{Mistele2024}. This relocates the split to a selection systematic rather
than the intrinsic type-dependence the law forbids; but it rests on a smaller, noisier sample and does not erase the
released-catalogue contrast. We present it as an open selection-versus-intrinsic question, decidable by the deeper,
better-isolated samples of Euclid and Rubin this decade.

\section{What lensing does and does not establish}
The honest ledger, both ways. \emph{Resolved:} the no-slip identity is real and removes the logical independence of
the Bullet, cluster strong$+$weak lensing, the bulk galaxy-lensing RAR, and Einstein radii from the dark-matter case;
the galaxy--galaxy-lensing RAR passes at the framework's own $a_0$ (a soft pass, inside the data
scatter~\cite{Brouwer2021}), curing the historical under-lensing of non-relativistic MOND. \emph{Conceded, genuinely
independent:} the cluster residual $\eta\!\simeq\!2.3$ (real, confirmed by lensing independently of hydrostatic bias,
shared with all MOND, and uncured here---the AeST scalar mass term is a candidate but is wrong-signed at clusters and
squeezed by galaxy lensing~\cite{Durakovic2024,Mistele2023}); and the CMB third peak. \emph{Contested:} the morphology
split. \emph{Open:} lensed-quasar substructure. Lensing is therefore not, in this framework, an independent detector
of a dark particle; it is a second window on the one shared cluster residual, with the framework's largest specifically
lensing tension---the morphology split---now in dispute.

\section{Conclusion}
``Gravitational lensing proves dark matter'' is, for a relativistic modified-dynamics theory with no gravitational
slip, a statement that double-counts. The lensing mass equals the dynamical mass by construction, so the
deflection-of-light arguments restate the cluster dynamical residual rather than corroborate it independently. This
does not make the cluster residual disappear---it is real, and conceded---nor does it touch the genuinely independent
constraints from the cosmic-microwave-background power spectrum. But it removes lensing from the list of independent
proofs of a dark particle and returns the dark-sector question to where the framework already locates it: a single,
soft, shared residual at the centres of clusters, and the one falsifiable thing the framework alone predicts, a
redshift-declining $a_0(z)\propto\sqrt{\rho_{\rm DE}(z)}$.

\begin{acknowledgments}
The acceleration-scale footing audit, the eRASS1 cluster residual at the framework's own $a_0$, and the morphology-split
cancellation were performed and cross-checked numerically; the no-slip identity is quoted from the primary AeST
literature. Both-ways audit ledgers are archived with this work.
\end{acknowledgments}

\begin{thebibliography}{99}
\bibitem{Clowe2006} D.~Clowe \emph{et al.}, Astrophys.\ J.\ \textbf{648}, L109 (2006) [Bullet cluster].
\bibitem{Milgrom1983} M.~Milgrom, Astrophys.\ J.\ \textbf{270}, 365 (1983).
\bibitem{Brouwer2021} M.~Brouwer \emph{et al.}, Astron.\ Astrophys.\ \textbf{650}, A113 (2021) [KiDS lensing RAR].
\bibitem{Skordis2021} C.~Skordis, T.~Z\l o\'snik, Phys.\ Rev.\ Lett.\ \textbf{127}, 161302 (2021) [AeST].
\bibitem{Verwayen2024} P.~Verwayen, C.~Skordis, C.~Z\l o\'snik, Mon.\ Not.\ R.\ Astron.\ Soc.\ \textbf{531}, 272 (2024)
[AeST quasistatic solutions; no slip].
\bibitem{Mistele2023} T.~Mistele, S.~McGaugh, S.~Hossenfelder, Astron.\ Astrophys.\ \textbf{676}, A100 (2023)
[AeST galaxy lensing].
\bibitem{Bulbul2024} E.~Bulbul \emph{et al.}, Astron.\ Astrophys.\ \textbf{685}, A106 (2024) [eRASS1 clusters].
\bibitem{Cha2026} S.~Cha \emph{et al.}, Astrophys.\ J.\ (2026), arXiv:2604.10811 [consistent-QUMOND Bullet lensing].
\bibitem{JWSTbullet} arXiv:2512.03150 (2025) [JWST joint-lensing Bullet, minor-merger reclassification].
\bibitem{Mistele2024} T.~Mistele \emph{et al.}, J.\ Cosmol.\ Astropart.\ Phys.\ \textbf{04}, 020 (2024)
[lensing RAR, isolation cut].
\bibitem{Durakovic2024} A.~Durakovi\'c, C.~Skordis (2024) [AeST scalar mass term in clusters].
\end{thebibliography}

\end{document}
